% sections/05-conclusion.tex

\section{Conclusion}
\label{sec:conclusion}

We have presented the first formal privacy analysis of the Kulldorff spatial scan
statistic, addressing a gap at the intersection of spatial epidemiology and differential
privacy.
Three theorems establish a complete pipeline: a sensitivity bound that quantifies
how much one reporting node can shift the log-likelihood ratio (Theorem~\ref{thm:sensitivity}),
a private window selection mechanism based on density-normalized utility
(Theorem~\ref{thm:mechanism}), and a composition analysis that governs the system's
operational lifetime under prospective (weekly, repeated) releases
(Theorem~\ref{thm:composition}).

The key insight is that $k_w$ (the number of reporting nodes inside a candidate
window) does not appear in the LLR formula because it is irrelevant to outbreak
detection, but it is the central parameter for privacy: it determines how much any
single participant can influence the statistic.
The density-normalized utility function $u(w) = \LLR(w)/\Delta(k_w)$ bridges the
two domains, enabling a clean application of the exponential mechanism with global
sensitivity~$1$.

\paragraph{Limitations.}
Theorem~\ref{thm:sensitivity} currently provides a candidate form for $f(\kw)$;
the formal proof and tightness characterization are the primary open problem.
The thesis is scoped to the Poisson LLR and primary cluster selection only;
the Bernoulli variant and multiple cluster privatization are future work.

\paragraph{Future work.}
\begin{itemize}
    \item Formally prove and tighten Theorem~\ref{thm:sensitivity}.
    \item Extend to multiple cluster selection using the sparse vector technique~\cite{dwork2013}.
    \item Apply to the VISHC wastewater surveillance data once available.
    \item Investigate smooth sensitivity~\cite{nissim2007} as an alternative to
          global sensitivity for further utility improvement.
\end{itemize}
